\section{Sujet n\textsuperscript{o}\,3}
\noindent\begin{minipage}{0.5\linewidth}\footnotesize Office du Baccalauréat du Cameroun\end{minipage}%
\hfill\begin{minipage}{0.45\linewidth}\footnotesize\raggedleft Examen : \textbf{Baccalauréat} · Série : \textbf{C/E}\\
Épreuve : \textbf{Mathématiques} · Durée : \textbf{4\,h} · Coef. : \textbf{7}\end{minipage}
\par\smallskip{\color{onbo}\rule{\linewidth}{1pt}}

\subsection*{Partie A — Évaluation des ressources \hfill (15 points)}

\noindent\textbf{Exercice 1.} \hfill\textit{(3 points)}\\
Le tableau suivant donne le nombre de sacs de ciment $x_i$ (en centaines) livrés sur un chantier
et le coût total de transport $y_i$ (en milliers de francs) pour cinq livraisons :
\par\smallskip
\begin{center}\footnotesize
\begin{tabular}{|c|c|c|c|c|c|}\hline
$x_i$ & $1$ & $2$ & $3$ & $4$ & $5$\\\hline
$y_i$ & $30$ & $50$ & $55$ & $70$ & $95$\\\hline
\end{tabular}
\end{center}
\begin{enumerate}[label=\arabic*)]
\item Calculer les moyennes $\bar x$ et $\bar y$, puis la variance $V(x)$ et la covariance $\mathrm{cov}(x,y)$. \hfill\textit{(1,5 pt)}
\item Déterminer une équation de la droite de régression de $y$ en $x$ par la méthode des moindres carrés. \hfill\textit{(1 pt)}
\item Estimer le coût de transport pour une livraison de $600$ sacs de ciment. \hfill\textit{(0,5 pt)}
\end{enumerate}

\smallskip
\noindent\textbf{Exercice 2.} \hfill\textit{(3 points)}\\
\begin{enumerate}[label=\arabic*)]
\item Déterminer le PGCD de $1\,365$ et $819$ à l'aide de l'algorithme d'Euclide. \hfill\textit{(1 pt)}
\item En déduire, à l'aide de l'algorithme d'Euclide, un couple $(u,v)\in\Z^2$ tel que
$1\,365\,u+819\,v=\mathrm{PGCD}(1365,819)$. \hfill\textit{(1,25 pt)}
\item Déterminer le reste de la division euclidienne de $3^{2026}$ par $7$. \hfill\textit{(0,75 pt)}
\end{enumerate}

\smallskip
\noindent\textbf{Exercice 3.} \hfill\textit{(4 points)}\\
Soit $f$ la fonction définie sur $\R$ par $f(x)=(2-x)\mathrm{e}^{x}$, de courbe $(\mathcal{C})$.
\begin{enumerate}[label=\arabic*)]
\item Déterminer les limites de $f$ en $-\infty$ et en $+\infty$. \hfill\textit{(0,75 pt)}
\item Montrer que $f'(x)=(1-x)\mathrm{e}^{x}$ et dresser le tableau de variations de $f$. \hfill\textit{(1,25 pt)}
\item Vérifier que $f$ est solution de l'équation différentielle $y'-y=-\mathrm{e}^{x}$ ;
résoudre l'équation différentielle $y'-y=0$ sur $\R$. \hfill\textit{(1 pt)}
\item À l'aide d'une intégration par parties, calculer $\displaystyle J=\int_0^{1}(2-x)\mathrm{e}^{x}\dd x$. \hfill\textit{(1 pt)}
\end{enumerate}

\smallskip
\noindent\textbf{Exercice 4.} \hfill\textit{(5 points)}\\
L'espace est muni d'un repère orthonormé direct $(O\,;\,\vec\imath,\vec\jmath,\vec k)$. On considère
les points $A(1\,;0\,;2)$, $B(2\,;1\,;0)$ et $C(0\,;1\,;1)$.
\begin{enumerate}[label=\arabic*)]
\item Calculer les coordonnées du vecteur $\vec{AB}\wedge\vec{AC}$ (produit vectoriel). \hfill\textit{(1,25 pt)}
\item En déduire l'aire du triangle $ABC$. \hfill\textit{(1 pt)}
\item Déterminer une équation cartésienne du plan $(ABC)$. \hfill\textit{(1,25 pt)}
\item Calculer la distance du point $D(3\,;3\,;3)$ au plan $(ABC)$, puis le volume du tétraèdre $ABCD$. \hfill\textit{(1,5 pt)}
\end{enumerate}

\subsection*{Partie B — Évaluation des compétences \hfill (5 points)}

\noindent\textit{L'entreprise BTP « Bâtir Cameroun » réalise la dalle d'un immeuble à Yaoundé et
te consulte, en tant qu'élève de Terminale~C. Le chef de chantier veut d'abord aménager un dépôt
\emph{rectangulaire} contre un long mur déjà bâti (le côté le long du mur n'est pas clôturé) ;
il dispose de \SI{120}{\metre} de barrière. Ensuite, il coule le béton : la première semaine
$\SI{40}{\cubic\metre}$ sont coulés, et le volume hebdomadaire augmente de \SI{10}{\percent}
chaque semaine ; le chantier nécessite au total $\SI{500}{\cubic\metre}$ de béton. Enfin, à la
réception des parpaings, \SI{3}{\percent} sont fissurés ; un contrôleur en examine $6$ au hasard.}

\begin{enumerate}[label=\textbf{Tâche \arabic*.},leftmargin=2.4em]
\item Déterminer les dimensions du dépôt rectangulaire d'\emph{aire maximale} réalisable avec la
barrière, ainsi que cette aire. \hfill\textit{(1,5 pt)}
\item Déterminer le nombre de semaines nécessaires pour couler au moins les $\SI{500}{\cubic\metre}$ de béton. \hfill\textit{(1,5 pt)}
\item Déterminer la probabilité qu'\emph{aucun} des $6$ parpaings examinés ne soit fissuré. \hfill\textit{(1,5 pt)}
\end{enumerate}
\par\smallskip{\footnotesize\itshape Présentation générale : 0,5 point.}

% ════════════════════════════════════════════════════
\corrige{du sujet 3}

\noindent\textbf{Partie A — Exercice 1.}
1) $\bar x=\frac{1+2+3+4+5}{5}=3$ ; $\bar y=\frac{30+50+55+70+95}{5}=\frac{300}{5}=60$.
$V(x)=\frac15\sum x_i^2-\bar x^2=\frac{1+4+9+16+25}{5}-9=\frac{55}{5}-9=11-9=2$.
$\mathrm{cov}(x,y)=\frac15\sum x_iy_i-\bar x\,\bar y$. Or $\sum x_iy_i=30+100+165+280+475=1050$,
donc $\mathrm{cov}(x,y)=\frac{1050}{5}-180=210-180=30$.
2) Pente $a=\frac{\mathrm{cov}(x,y)}{V(x)}=\frac{30}{2}=15$ ; $b=\bar y-a\bar x=60-15\times3=15$.
Droite : $y=15x+15$.
3) $600$ sacs $=6$ centaines, donc $x=6$ : $y=15\times6+15=105$, soit \SI{105000}{F}.

\noindent\textbf{Exercice 2.}
1) $1365=1\times819+546$ ; $819=1\times546+273$ ; $546=2\times273+0$. PGCD $=273$.
2) On remonte : $273=819-1\times546=819-(1365-819)=2\times819-1365$. Donc
$1365\times(-1)+819\times2=273$ : un couple est $(u,v)=(-1\,;2)$.
3) Modulo $7$ : $3^1\equiv3,\ 3^2\equiv2,\ 3^3\equiv6,\ 3^4\equiv4,\ 3^5\equiv5,\ 3^6\equiv1\ [7]$
(ordre $6$). Comme $2026=6\times337+4$, $3^{2026}\equiv3^4\equiv4\ [7]$. Le reste est $\mathbf 4$.

\noindent\textbf{Exercice 3.}
1) En $-\infty$ : $\mathrm e^x\to0$ et par croissances comparées $x\,\mathrm e^x\to0$, donc $f(x)=(2-x)\mathrm e^x\to0^+$.
En $+\infty$ : $2-x\to-\infty$ et $\mathrm e^x\to+\infty$, donc $f(x)\to-\infty$.
2) $f'(x)=(-1)\mathrm e^x+(2-x)\mathrm e^x=(1-x)\mathrm e^x$. Signe de $1-x$ : $f'>0$ sur $]-\infty;1[$,
$f'<0$ sur $]1;+\infty[$. Maximum en $x=1$ : $f(1)=(2-1)\mathrm e=\mathrm e$. $f$ croît puis décroît.
3) $f'(x)-f(x)=(1-x)\mathrm e^x-(2-x)\mathrm e^x=\big((1-x)-(2-x)\big)\mathrm e^x=-\mathrm e^x$ : $f$ est bien solution.
L'équation $y'-y=0$, soit $y'=y$, a pour solutions $y(x)=K\mathrm e^{x}$, $K\in\R$.
4) IPP avec $u=2-x$, $v'=\mathrm e^x$, $u'=-1$, $v=\mathrm e^x$ :
$J=\big[(2-x)\mathrm e^x\big]_0^1+\int_0^1\mathrm e^x\dd x=\big((2-1)\mathrm e-(2)\cdot1\big)+\big[\mathrm e^x\big]_0^1
=(\mathrm e-2)+(\mathrm e-1)=2\mathrm e-3$.

\noindent\textbf{Exercice 4.}
1) $\vec{AB}=(1\,;1\,;-2)$, $\vec{AC}=(-1\,;1\,;-1)$.
$\vec{AB}\wedge\vec{AC}=\big(1\cdot(-1)-(-2)\cdot1\,;\ (-2)\cdot(-1)-1\cdot(-1)\,;\ 1\cdot1-1\cdot(-1)\big)=(1\,;3\,;2)$.
2) Aire $=\frac12\|\vec{AB}\wedge\vec{AC}\|=\frac12\sqrt{1^2+3^2+2^2}=\frac12\sqrt{14}$.
3) Le plan $(ABC)$ a pour vecteur normal $\vec n=(1\,;3\,;2)$ : $x+3y+2z=d$. Avec $A(1;0;2)$ :
$d=1+0+4=5$. Équation : $x+3y+2z-5=0$.
4) $d(D,(ABC))=\frac{|3+3\times3+2\times3-5|}{\sqrt{14}}=\frac{|3+9+6-5|}{\sqrt{14}}=\frac{13}{\sqrt{14}}$.
Volume $=\frac13\times\text{aire}\times d=\frac13\times\frac{\sqrt{14}}2\times\frac{13}{\sqrt{14}}=\frac{13}{6}$.

\smallskip
\noindent\textbf{Partie B — Situation.}
\emph{Tâche 1.} Soit $x$ la largeur (deux côtés perpendiculaires au mur) et $y$ la longueur le long
du mur. La barrière couvre $2x+y=120$, donc $y=120-2x$. Aire $A(x)=x(120-2x)=120x-2x^2$.
$A'(x)=120-4x=0$ en $x=30$ ; $A''<0$ : maximum. Dimensions $x=\SI{30}{\metre}$, $y=120-60=\SI{60}{\metre}$ ;
aire maximale $A=30\times60=\SI{1800}{\metre\squared}$.

\emph{Tâche 2.} Volume coulé la semaine $n$ : $w_n=40\times1{,}1^{\,n-1}$ ($\SI{}{\cubic\metre}$).
Volume cumulé sur $n$ semaines : $T_n=40\cdot\frac{1{,}1^{\,n}-1}{1{,}1-1}=400\,(1{,}1^{\,n}-1)$.
On résout $T_n\ge500$ : $1{,}1^{\,n}-1\ge1{,}25$, soit $1{,}1^{\,n}\ge2{,}25$, d'où
$n\ge\frac{\ln2{,}25}{\ln1{,}1}\approx\frac{0{,}8109}{0{,}0953}\approx8{,}51$ : il faut \textbf{9 semaines}
($T_9=400(1{,}1^9-1)\approx400\times1{,}358=\SI{543}{\cubic\metre}$, alors que $T_8\approx\SI{457}{\cubic\metre}<500$).

\emph{Tâche 3.} Chaque parpaing est fissuré avec probabilité $p=0{,}03$ ; $X\sim\mathcal B(6\,;0{,}03)$.
$P(X=0)=(1-0{,}03)^6=(0{,}97)^6\approx\mathbf{0{,}833}$, soit environ \SI{83,3}{\percent}.
